\documentclass[10pt,conference]{IEEEtran}
\usepackage[T1]{fontenc}
\usepackage{mathptmx}
\usepackage{amsmath}
\usepackage{booktabs}
\usepackage{graphicx}
\usepackage{xcolor}
\usepackage[hidelinks]{hyperref}

\newcommand{\tbd}[1]{\textcolor{red}{#1}}

\begin{document}
\title{Predicting Heart Disease from Survey Answers:\\
the Codebook Matters More than the Regularizer}

\author{
  \IEEEauthorblockN{Youssef Bensalem, Youssef Masmoudi, Ilias Sokolis}
  \IEEEauthorblockA{CS-433 Machine Learning, EPFL, Fall 2026}
}

\maketitle

\begin{abstract}
We predict coronary heart disease (MICHD) for 328,135 respondents of the 2015
BRFSS telephone survey from their 321 self-reported answers, using logistic
regression trained by gradient descent and scored by F1. Cross-validation shows
that the model cannot overfit ($N/D \approx 2{,}100$), so L2 regularization only
hurts, and the decisive choices are the decision threshold and how answers are
encoded. BRFSS stores ``don't know'' and ``refused'' as per-variable numeric
codes that a linear model reads as extreme values; recoding them to missing with
the official codebook is the single representation change that helps
($+0.013$ CV F1). Our final model reaches a 5-fold CV F1 of 0.425.
We also report three negative results: one-hot encoding, squared features and
class rebalancing.
\end{abstract}

\section{Introduction}
The task is binary classification with a strong imbalance: 8.83\% of training
rows are positive. Accuracy is therefore useless (predicting ``healthy'' for
everyone scores 91\%), and the competition metric is the F1 score of the
positive class. Predicting \emph{everyone} positive already gives
F1~$=2p/(1+p)=0.162$, which is the floor any model has to beat. Every feature is
a questionnaire answer: there are no lab values, and 44.7\% of all cells are
missing, mostly because of the survey's skip logic.

Only NumPy is allowed, so we stay with linear models and spend the effort on
validation and on how the answers are represented. All full-data results are
reproduced by \texttt{run.py} and \texttt{ablation.py} in our repository.

\section{Preprocessing}
\label{sec:prep}
Every statistic below is fitted inside each training fold only.

\textbf{1. Sentinel codes $\to$ missing.} BRFSS codes a non-answer just above each
variable's valid range, and the code differs per variable: 7/9 for
\texttt{GENHLTH} (scale 1--5), 77/99 for the 0--30 day counts (where 7 and 9 are
real answers), 7777/9999 for weight in pounds, 3 for \texttt{\_AGE65YR}, 14 for
\texttt{\_AGEG5YR}. Left in place, these codes are read as magnitudes.
Fitting logistic regression on \texttt{GENHLTH} alone gives a slope of $+0.80$
per standard deviation, i.e.\ the odds are multiplied by 2.2 per step. A
refusal (9) standardizes to $+5.8\sigma$, 3.6$\sigma$ beyond the worst real
answer, where the model predicts an 88\% risk against an observed 14\%. The
observed risk rises from 1 to 5 and then \emph{falls} at 7 and 9, a shape no
single coefficient can fit. We took the codes of 242 of the 321 columns from the
official CDC 2015 SAS format files~\cite{brfss} (\texttt{brfss\_codebook.py}).
We first tried to infer them from the data (``a code is a sentinel if the value
below it never occurs''), but this rule cannot find \texttt{\_AGE65YR}$=3$ or
\texttt{\_AGEG5YR}$=14$, which follow their real values without a gap, and a
cruder ``7 and 9 are sentinels'' rule would have erased real day counts.

\textbf{2. Column filtering.} We drop columns that are more than 50\% missing
after step~1, constant columns, and 28 survey administration columns (record
numbers, sampling weights, interview dates, telephone screening, and the state
FIPS code, whose numeric order is meaningless). 152 columns remain.

\textbf{3. Standardization, then missing $\to 0$.} Each column is centred and
scaled with the mean and standard deviation of its \emph{observed} values;
missing entries are then set to 0. After centring, 0 is the column mean, so a
missing answer adds nothing to the score $w^\top z$. Any other value is a claim
about where the respondent sits on the health axis. Imputing a raw 0 before
scaling would be wrong, because 0 lies outside the valid range of most answers.

\textbf{4. Two counts and an intercept.} Recoding a refusal to missing merges
it with ``never asked''. To keep the distinction we add, per row, the
standardized number of refused answers and of questions never asked. Rows
with at least three refusals are 10.4\% positive, compared with 4.7\% for rows
with none. Together with the intercept this gives $D=155$.

\section{Models and Validation}
\textbf{Methods.} We implemented the six required methods (MSE by gradient
descent and SGD, least squares and ridge by normal equations, logistic and
L2-regularized logistic regression by gradient descent), all checked against
the course's reference values. Least squares on $\{0,1\}$ labels serves as
a baseline. The final model is logistic regression trained by full-batch
gradient descent from $w=0$.

\textbf{Step size and iterations.} After 500 steps the training loss is 0.2292
for $\gamma=0.1$ and 0.2258 for $\gamma=1$, and gradient descent diverges for
$\gamma\ge 5$ (loss 0.697, above the 0.693 of $w=0$). We use $\gamma=1$ and 500
iterations. On a held-out 20\%, 4000 iterations change validation F1 by only
$+0.001$, and validation loss is lowest at 500.

\textbf{Cross-validation.} We use 5 folds (seed 1). Preprocessing is refitted
inside every training fold, so no validation row influences it; all five folds
keep exactly the same columns as a fit on the full data. The decision
threshold plays no part in training, so we cache each fold's validation
probabilities and sweep 56 thresholds from 0.05 to 0.60 in steps of 0.01 at no
extra cost. We select the $(\lambda,\text{threshold})$ pair with the best mean
validation F1, then refit preprocessing and weights on all training rows. The
CV estimate is slightly pessimistic, because each fold trains on 80\% of the
data. It is also a good predictor of the leaderboard: an earlier configuration
scored 0.393 in CV (on a 10\% subsample) and 0.403 on AIcrowd.

\section{Results}
Table~\ref{tab:ablation} adds one change at a time, on the full training set
with the same five folds. Figure~\ref{fig:thr} shows F1 as a function of the
threshold.

\begin{table}[t]
\centering
\caption{Ablation on the full training set, 5-fold CV, mean $\pm$ std over
folds. Each logistic row adds one change to the row above. $\bar p$ is the mean
predicted probability; the base rate is 0.0883.}
\label{tab:ablation}
\footnotesize
\setlength{\tabcolsep}{3pt}
\begin{tabular}{lcccc}
\toprule
Configuration & $D$ & CV F1 & thr. & $\bar p$ \\
\midrule
Predict all positive & -- & 0.162 & -- & -- \\
Least squares, final features & 155 & $0.416\pm0.004$ & 0.21 & -- \\
\midrule
Logistic, $\gamma{=}0.1$, $\lambda{=}1$, all columns & 173 & $0.380\pm0.004$ & 0.50 & 0.455 \\
+ $\lambda=0$ & 173 & $0.408\pm0.003$ & 0.21 & 0.098 \\
+ $\gamma=1$ & 173 & $0.412\pm0.003$ & 0.20 & 0.088 \\
+ drop survey metadata & 154 & $0.412\pm0.003$ & 0.20 & 0.088 \\
+ sentinels $\to$ missing & 153 & $0.425\pm0.005$ & 0.20 & 0.088 \\
+ refused / never-asked counts & 155 & $\mathbf{0.425\pm0.005}$ & 0.20 & 0.088 \\
\midrule
Final model, threshold 0.5 & 155 & $0.226\pm0.005$ & 0.50 & 0.088 \\
\bottomrule
\end{tabular}
\end{table}

\begin{figure}[t]
\centering
\includegraphics[width=\columnwidth]{threshold_f1.pdf}
\caption{5-fold CV F1 against the decision threshold (full training set).}
\label{fig:thr}
\end{figure}

\textbf{The threshold is the largest single lever.} The final model scores
0.226 at the default threshold of 0.5 and 0.425 at 0.20. For a calibrated
classifier, the F1-optimal threshold is half the achievable F1~\cite{lipton},
here $0.425/2 \approx 0.21$, which is close to the 0.20 we select.

\textbf{Regularization does not help.} With $N/D\approx 2{,}100$ there is
almost no variance for L2 to control, and the CV and leaderboard scores agree,
so there is no overfitting to fix. Removing $\lambda=1$ is worth $+0.028$
(Table~\ref{tab:ablation}). Shrinking $w$, intercept included, pulls every
probability towards $\sigma(0)=0.5$: the regularized model has $\bar p=0.455$
and its best threshold is 0.5. This explains why our earlier submission
needed a threshold of 0.4. At $\lambda=0$, $\bar p$ equals the base rate, so
the model is calibrated. The grid in \texttt{run.py} gives CV F1
0.4248, 0.4252 and 0.4208 for $\lambda = 0$, 0.001 and 0.01. The selection rule
picks $\lambda=0.001$, but its lead is ten times smaller than the fold standard
deviation, so in practice $\lambda=0$ and $\lambda=0.001$ are equivalent. Larger values could not be tested at $\gamma=1$: each step
multiplies $w$ by $(1-2\gamma\lambda)$, so $\lambda=1$ gives the update
$w\leftarrow -w-\nabla L$, and $\lambda=0.1$ already oscillates (CV F1 0.16).
At $\gamma=0.1$ on a 10\% subsample, F1 fell steadily from 0.402
($\lambda=0$) to 0.393 ($\lambda=0.1$) and 0.341 ($\lambda=5$).
Increasing $\gamma$ from 0.1 to 1 adds $+0.004$, because 500 steps of
$\gamma=0.1$ have not converged ($\bar p=0.098$).


\textbf{Sentinel recoding is the representation change that helps.}
Dropping the metadata columns is neutral, which we expected: it is a
correctness fix, since record numbers and sampling weights carry no
information about the respondent. Recoding the sentinels adds $+0.013$, about
three times the fold standard deviation.
How the codes are neutralized matters. Excluding them from the mean and
standard deviation while still passing the raw value is catastrophic (CV F1
0.22 on a 10\% subsample, against 0.40 when they are left untouched): the
codes had been inflating $\sigma$ and so compressing themselves, and without
them they land even further out. Adding one ``refused'' indicator per column
(264 columns, on the subsample) gains nothing, and neither do the two counts
($+0.00002$). The association between refusals and the outcome is real, but it
is already explained by covariates such as age and general health.

\textbf{Negative results} (10\% subsample). \emph{(i) One-hot encoding}
the 97 kept columns that are categorical codes with 3--10 levels (772--1,062
columns) \emph{lowered} CV F1 to 0.20--0.27, even after 3000 iterations. The
cause is conditioning, not overfitting: a standardized dummy for a level held
by 0.1\% of rows takes the values $-0.03$ and $+31.6$, and gradient descent
with one step size cannot converge along such directions. The best threshold
sat at the bottom of the sweep and $\bar p$ far below the base rate.
\emph{(ii) Squared features} failed for the same reason (0.32--0.37).
\emph{(iii) Undersampling negatives} to 1:1, 2:1, 4:1 and 8:1 gave 0.355,
0.374, 0.392 and 0.402, against 0.402 without it, while the best threshold
moved from 0.70 to 0.24. Rebalancing a calibrated model only shifts its
optimal threshold~\cite{elkan}.

\section{Discussion}
Our final model ($\lambda=0.001$, threshold 0.20) reaches a 5-fold CV F1 of
0.425. Its submission
(\texttt{run.py}) scores \tbd{0.4xx} on AIcrowd, up from 0.403 for our
previous best ($\lambda=1$, threshold 0.4, no sentinel recoding).
Since there is no overfitting, the remaining gap to the best teams
($\approx0.45$) has to come from representation, not regularization. The clearest
candidate is one-hot encoding restricted to levels with at least 1\% support,
which keeps the problem well conditioned, or a second-order optimizer that is
insensitive to conditioning. Two limitations apply. Both the features and the
label are self-reported, so the model predicts a \emph{reported} diagnosis.
Hyperparameters were also chosen by CV with the leaderboard as the only
external check, so we kept the number of submissions small.

\begin{thebibliography}{1}
\bibitem{brfss} Centers for Disease Control and Prevention, ``Behavioral Risk
Factor Surveillance System, 2015 codebook and SAS format files
(\texttt{Formas15.sas}, \texttt{Format15.sas}),''
\url{https://www.cdc.gov/brfss/annual_data/annual_2015.html}.
\bibitem{lipton} Z.~C. Lipton, C.~Elkan, and B.~Narayanaswamy, ``Optimal
thresholding of classifiers to maximize F1 measure,'' in \emph{ECML PKDD},
2014, pp.~225--239.
\bibitem{elkan} C.~Elkan, ``The foundations of cost-sensitive learning,'' in
\emph{IJCAI}, 2001, pp.~973--978.
\end{thebibliography}

\end{document}
